\scanpage{111}
% L1--2
\unsure{Examinons|Examiner} avec plus de soin les parties \unsure{antitypiques|antypiques|atypiques|analytiques} \inserted{et loc. constructibles} de $\unsure{x|X}$, quand on ne suppose pas $X$ qu. cpt. \ldots\ldots

\begin{center}\rule{1.3cm}{0.4pt}\\[-3pt]\rule{.85cm}{0.4pt}\end{center}

% L3
\uline{Th} Soit $X$ un schéma affine. Alors $X^{\mathrm{cons}}$ est compact

% L4--7
\uline{Corollaire} 1 Soit $X$ un préschéma \deleted{\unsure{Raisonnement|Précisément}}. \inserted{qu. séparé.} Alors $X$ est $T_0$, admet une base formée d'ouverts rétrocompacts et compacts pour la topologie constructible. En particulier, $X^{\mathrm{cons}}$ est loc. compact.

% L8--9
\uline{Corollaire \unsure{3|2}} Supposons \deleted{de plus} $X$ [quasi-compact et quasi-séparé] \deleted{Alors} quasi-noethérien. Alors

% L10--11
(i) $X^{\mathrm{cons}}$ est compact. (ii) Les parties constructibles de $X$ sont les parties ouvertes et fermées de $X^{\mathrm{cons}}$.

% L12--17
(iii) Les parties \deleted{\unsure{affines de|affines}} antitypiques de $X$ sont les parties \inserted{$E$} telles qu'il existe un morphisme qu. cpt $f:X'\to X$, d'image $E$. On peut prendre \inserted{pour} \deleted{$\unsure{E|X}$} $X'\to X$ un \uline{monomorphisme} \deleted{[\unsure{N.B. ce sont}} qu. cpt \unsure{au|--} (iv) Si $E$ et $X-E$ ont cette propriété, ils sont constructibles \ldots
% archive sheet number: 110
